
\subsubsection{2.1 Generalization Gap Studies}

Recht et al. (2019) revealed that ImageNet classifiers exhibit systematic accuracy drops of 11-15\% on ImageNet-V2, a reproduced test set following the original data collection methodology. Follow-up work found 3-5\% gaps on CIFAR-10.2, while Barbu et al. (2019) documented 40-45\% drops on ObjectNet, which tests object recognition under varied viewpoints and backgrounds. In NLP, McCoy et al. (2019) introduced HANS, revealing that BERT's 84\% MNLI accuracy collapses when syntactic heuristics are isolated.

These studies share a limitation: they measure gaps after constructing held-out test sets. They provide no method to predict which benchmarks will exhibit large gaps using only historical SOTA records.

\subsubsection{2.2 Benchmark Analysis}

PapersWithCode (2020) provides SOTA tracking for benchmarks, enabling quantitative analysis of benchmark dynamics. Thompson et al. (2020) analyzed compute scaling trends but did not address saturation measurement. Bowman and Dahl (2021) critiqued NLP benchmark culture, arguing that leaderboard climbing incentivizes overfitting to test set idiosyncrasies, but without quantitative saturation metrics.

\subsubsection{2.3 Saturation Detection}

The concept of benchmark saturation appears informally in ML discourse. Ethayarajh et al. (2022) measured dataset difficulty via V-usable information, while Rodriguez et al. (2021) analyzed annotation difficulty. Neither addresses saturation via improvement entropy.

DNSI differs by measuring the entropy of improvement patterns in SOTA histories rather than dataset properties directly.
